\subsection{The panel layer}
\label{sec:molecular-algebraic-induction-panel}

Katoh--Tanigawa's \emph{panel-hinge} framework is a
\emph{hinge-coplanar} body-hinge framework: at each body all incident
hinges lie in a common hyperplane \cite[p.\,647]{katohTanigawa2011}. The
molecular conjecture's content is that this coplanarity can be met
without losing rigidity; the Phase-18 body-hinge framework
(\cref{def:hinge-constraint}) carries free hinges with no coplanarity,
so the realization-existence nodes below need a \emph{panel layer} on
top. We take the panel-data form: a panel realization assigns each body
$v$ a hyperplane normal $n_v \in \R^{k+2}$, and the hinge at an edge
$e = uv$ is the codimension-2 intersection $\mathrm{panel}(u) \cap
\mathrm{panel}(v)$. Its supporting $k$-extensor --- the screw-space
element of \cref{def:hinge-constraint} the rigidity matrix constrains ---
is the Grassmann--Cayley meet of the two panels, equivalently the
complement iso of the join of the two normals. Coplanarity (both hinges
at $v$ lie in $\mathrm{panel}(v)$ by construction) and transversality
both live in the extensor algebra, so no affine-subspace-intersection
plumbing is needed.

\begin{definition}[Panel support extensor]
  \label{def:panel-support-extensor}
  \lean{CombinatorialRigidity.Molecular.panelSupportExtensor,
        CombinatorialRigidity.Molecular.normalsJoin,
        CombinatorialRigidity.Molecular.panelSupportExtensor_ne_zero_iff}
  \leanok
  \uses{def:join, def:meet-complement-iso, def:hinge-constraint}
  For two panel normals $n_1, n_2 \in \R^{k+2}$, the \emph{panel support
  extensor} is the supporting $k$-extensor $C(p(e)) \in \Lambda^k
  \R^{k+2}$ of the codimension-2 intersection of the two panels, given as
  the complement iso (\cref{def:meet-complement-iso}) of the grade-2 join
  $n_1 \wedge n_2 \in \Lambda^2 \R^{k+2}$ (\cref{def:join}). Since
  $k+2-2 = k$ is the screw-space grade, the result lands in the screw
  space $\Lambda^k \R^{k+2}$ of \cref{def:hinge-constraint} and feeds the
  hinge constraint verbatim. It is nonzero exactly when the two normals
  are linearly independent --- the algebraic form of the two panels
  meeting transversally in a genuine hinge --- the only general-position
  condition the panel layer needs.
\end{definition}

\begin{definition}[Panel-hinge framework and its body-hinge interpretation]
  \label{def:panel-hinge-framework}
  \lean{CombinatorialRigidity.Molecular.PanelHingeFramework,
        CombinatorialRigidity.Molecular.PanelHingeFramework.toBodyHinge,
        CombinatorialRigidity.Molecular.BodyHingeFramework.IsHingeCoplanar,
        CombinatorialRigidity.Molecular.BodyHingeFramework.isHingeCoplanar_toBodyHinge,
        CombinatorialRigidity.Molecular.PanelHingeFramework.IsGeneralPosition,
        CombinatorialRigidity.Molecular.PanelHingeFramework.supportExtensor_ne_zero_of_isGeneralPosition,
        CombinatorialRigidity.Molecular.PanelHingeFramework.isGeneralPosition_withMomentNormals,
        CombinatorialRigidity.Molecular.PanelHingeFramework.isGeneralPosition_ofParam}
  \leanok
  \uses{def:panel-support-extensor, def:hinge-constraint}
  A \emph{panel-hinge framework} is a multigraph $G = (V,E)$ together with a
  per-body panel normal $n_v \in \R^{k+2}$ (the pole of body $v$'s
  hyperplane $\mathrm{panel}(v)$). The hinge at an edge $e = uv$ is the
  codimension-2 intersection $\mathrm{panel}(u) \cap \mathrm{panel}(v)$,
  with supporting $k$-extensor the meet $C(p(e)) =
  \mathrm{panelSupportExtensor}(n_u, n_v)$
  (\cref{def:panel-support-extensor}). Its \emph{body-hinge interpretation}
  is the body-hinge framework (\cref{def:hinge-constraint}) on $G$ carrying
  exactly these supporting extensors, so the entire rigidity-matrix rank
  theory of \cref{sec:molecular-rigidity-matrix} applies verbatim. Every
  hinge incident to $v$ has $n_v$ as a join factor, hence lies in the single
  panel $\mathrm{panel}(v)$: the framework is \emph{hinge-coplanar} by
  construction. A body-hinge framework is hinge-coplanar
  ($\mathrm{IsHingeCoplanar}$) when it arises this way; this is the
  molecular-model constraint Katoh--Tanigawa's panel-hinge frameworks carry
  beyond the free-hinge body-hinge model, and the conjecture's content is
  that it can be met without dropping rigidity. The normals are in
  \emph{general position} ($\mathrm{IsGeneralPosition}$) when any two normals
  at distinct bodies are linearly independent; then every hinge joining two
  distinct bodies has a nonzero supporting extensor
  ($\mathrm{supportExtensor\_ne\_zero\_of\_isGeneralPosition}$) --- the
  transversality the rigid-block forest hinges feed into the block-triangular
  gluing of \cref{lem:case-I}. An explicit, dimension-free general-position
  assignment exists for any number of bodies: setting each body's normal to a
  point $(1, t, t^2, \dots, t^{k+1})$ on the moment curve at a distinct
  parameter $t$ makes any two normals at distinct bodies independent (the
  $2 \times 2$ Vandermonde minor has determinant $t' - t \neq 0$), so an
  injective real assignment of parameters yields
  $\mathrm{IsGeneralPosition}$
  ($\mathrm{isGeneralPosition\_withMomentNormals}$, and from a bare graph and
  endpoint selector directly via $\mathrm{isGeneralPosition\_ofParam}$ --- the
  realization-side entry point for the Case-I assembly)
  --- where standard-basis normals would cover only $|V| \le k+2$.
\end{definition}

\subsection{The generic-rank reconciliation (Proposition 1.1, analytic half)}
\label{sec:molecular-algebraic-induction-prop11}

The matroidal half of Tay--Whiteley's Proposition~1.1 --- the bridge
$\operatorname{def}(\tilde G) = \operatorname{corank} M(\tilde G)$ ---
landed green in Phase~19 (\cref{thm:def-eq-corank}). The analytic half
relocated forward to this phase: it wires the rigidity-matrix rank
$R(G,p)$ of \cref{sec:molecular-rigidity-matrix} to the matroid corank,
and depends on the generic-rank argument (Claim~6.4) that drives the
whole induction.

\begin{proposition}[Generic rank equals $D(|V|-1) - \operatorname{def}$; KT Proposition 1.1, analytic half]
  \label{prop:rigidity-matrix-prop11}
  \lean{CombinatorialRigidity.Molecular.rigidityMatrix_prop11}
  \leanok
  \uses{def:rigidity-matrix, def:D-deficiency, thm:def-eq-corank, def:dof-generic, lem:genericity-device}
  For a generic panel-hinge realization $(G,p)$,
  \[
    \operatorname{rank} R(G,p) = D(|V|-1) - \operatorname{def}(\tilde G).
  \]
  Equivalently $G$ has a rigid panel-hinge realization iff
  $\tilde G = (D-1)\cdot G$ contains $D$ edge-disjoint spanning trees,
  reconciling the honest rank form of
  \cref{sec:molecular-rigidity-matrix} with the body-hinge existence
  form (\cref{thm:body-hinge-tay}) of Phase~16. This is
  Jackson--Jord\'an's Theorem~6.1, geometric side
  \cite{jacksonJordan2009}.
\end{proposition}
\begin{proof}
  \leanok
  \uses{thm:theorem-55, lem:trivial-motions-rank-bound,
        lem:motions-mono-of-graph-le}
  The deficiency is the corank of $M(\tilde G)$
  (\cref{thm:def-eq-corank}), so the claim is that a generic realization
  attains the maximum rank $D(|V|-1) - \operatorname{corank}$. The target
  equality is the realization hypothesis \cref{def:rank-hypothesis} at
  $k' = \operatorname{def}$; it is pinned by two inequalities. The
  genericity-free upper bound $\operatorname{rank} R(G,p) \le D(|V|-1) -
  \operatorname{def}$ is the codimension form
  \cref{lem:trivial-motions-rank-bound} together with the deficiency count: a
  partition $P$ attaining $\operatorname{def}(\tilde G)$ contributes
  $\operatorname{def}$ independent motions beyond the $D$ trivial ones --- this
  holds for \emph{every} realization, no genericity, and is discharged in-proof
  from the Phase-19 partition machinery. The reverse (generic max-rank) bound
  $\operatorname{rank} R(G,p) \ge D(|V|-1) - \operatorname{def}$ --- equivalently
  $\dim Z(G,p) \le D + \operatorname{def}$, carried as the hypothesis
  $\mathtt{hgen}$ --- holds at a generic realization: delete $G$ down to a
  minimal $k$-dof spanning subgraph, realize it at full generic rank by
  Theorem~5.5 (\cref{thm:theorem-55}, the Claim~6.4 generic-rank argument
  selecting the attaining point), then re-add the deleted edges, which only
  grows the rank (\cref{lem:motions-mono-of-graph-le}). The proposition's only
  further inputs are the dimension fixing $n = k+1$ and the genuine-hinge
  condition $C(e) \neq 0$ the partition cut needs. The construction supplying
  $\mathtt{hgen}$ at every deficiency is Theorem~5.6 at $d = 3$
  (\cref{thm:theorem-55-6-d3}). Applied to a spanning $G$ ($V(G) = V_\alpha$),
  the $V(G)$-relative motive of \cref{def:rank-hypothesis} coincides with the
  absolute $\dim Z = D + \operatorname{def}$.
\end{proof}

\subsection{The induction skeleton and base case}
\label{sec:molecular-algebraic-induction-base}

\begin{definition}[Generic-rank hypothesis (6.1), $V(G)$-relative --- rank-form substrate]
  \label{def:rank-hypothesis}
  \lean{CombinatorialRigidity.Molecular.BodyHingeFramework.IsInfinitesimallyRigidOn,
        CombinatorialRigidity.Molecular.PanelHingeFramework.HasFullRankRealization,
        CombinatorialRigidity.Molecular.PanelHingeFramework.HasGenericFullRankRealization,
        CombinatorialRigidity.Molecular.BodyHingeFramework.RankHypothesis}
  \leanok
  \uses{def:rigidity-matrix, def:k-dof}
  For a minimal $k$-dof-graph $G$ with $|V| \ge 2$, the
  \emph{realization hypothesis} $(6.1)$ asserts the existence of a
  panel-hinge realization $(G,p)$ with
  $\operatorname{rank} R(G,p) = D(|V|-1) - k$.

  \emph{The rank must be read $V(G)$-relative.} The rigidity rows of
  $R(G,p)$ come only from the edges of $G$, so $\operatorname{rank} R(G,p) =
  \dim \operatorname{span}(\text{rigidity rows})$ depends only on $V(G)$ ---
  it is blind to any ambient bodies not in $V(G)$. The target form the Lean
  carries is therefore the \emph{rank} equality
  \[
    \dim \operatorname{span}\bigl(\text{rigidity rows of } (G,p)\bigr)
      \;=\; D\,(|V(G)| - 1) - \operatorname{def}(\tilde{G}),
  \]
  equivalently ``every infinitesimal motion is constant on $V(G)$'' when
  $\operatorname{def} = 0$. This is \emph{not} the same as the null-space equality
  $\dim Z(G,p) = D + k$ over the whole screw-assignment space
  $\R^{D|V_\alpha|}$ of the ambient body type; the rank form above is the
  $\alpha$-independent motive the realization layer is stated against.

  \textbf{Lean forms.}  Two predicates pin this definition.

  $\mathtt{HasFullRankRealization}\,k\,G$ is the bare existential
  $\exists\,Q,\; Q.\mathrm{graph} = G \;\wedge\;
  Q.\mathrm{toBodyHinge}.\mathrm{IsInfinitesimallyRigidOn}\,V(G)$,
  used as the induction motive of Theorem~5.5.
  It is equivalent to the rank equality (B1,
  $\mathtt{isInfinitesimallyRigidOn\_vertexSet\_iff\_finrank\_span\_rigidityRows}$)
  when $V(G)$ is nonempty.  It remains \emph{weaker} than KT's (6.1) ---
  a degenerate welded framework satisfies it on every connected $0$-dof
  graph --- and is slated for deletion once the spine restates onto the
  honest genuine-hinge motive (\cref{def:genuine-hinge-realization}).

  $\mathtt{HasGenericFullRankRealization}\,k\,n\,G$ is the strengthened motive
  that carries the rank equality as an explicit conjunct (at deficiency $n$,
  alongside general position and algebraic independence of the normals);
  it is the induction substrate for the generic-motive reduction
  $\mathtt{theorem\_55\_generic}$.  The rank conjunct is the
  $\mathbb{Z}$-cast form
  $(\dim \operatorname{span}(\text{rigidity rows}) : \mathbb{Z}) =
    D\,(|V(G)| - 1) - \operatorname{def}(\tilde{G})$,
  matching the $\mathbb{Z}$ arithmetic of the genuine-hinge motive
  (\cref{def:genuine-hinge-realization}).
\end{definition}

\begin{definition}[Generic-rank hypothesis (6.1) --- genuine-hinge form; KT eq.\ (6.1)]
  \label{def:genuine-hinge-realization}
  \lean{CombinatorialRigidity.Molecular.HasPanelRealization,
    CombinatorialRigidity.Molecular.ExtensorInPanel,
    CombinatorialRigidity.Molecular.hasPanelRealization_of_generic}
  \leanok
  \uses{def:rank-hypothesis, def:panel-hinge-framework, def:D-deficiency}
  A graph $G$ \emph{has a genuine-hinge panel realization} at dimension $k$ and
  defect parameter $n$ when there exist a body-hinge framework $F$ on $G$ and a
  panel-normal assignment $\mathtt{normal} : V \to \mathbb{R}^{k+2}$ such that:
  every vertex has a nonzero normal; every link's supporting $k$-extensor is
  nonzero and lies in both endpoint panels (each endpoint's hyperplane
  $\mathtt{normal}(v)^\perp$ contains $k$ spanning points of the extensor,
  $\mathtt{ExtensorInPanel}$); and the rigidity-row span attains the target rank
  $D(|V(G)|-1) - \operatorname{def}(\tilde G)$ (KT eq.~(6.1), $\mathbb Z$-cast form).

  Every generic realization (\cref{def:panel-hinge-framework}) with
  $|V| \ge 2$ yields a genuine-hinge realization
  ($\mathtt{hasPanelRealization\_of\_generic}$): panel normals in general position
  give nonzero extensors (general position forces transversality on each link);
  the meet-decomposition lemma furnishes the $\mathtt{ExtensorInPanel}$ witnesses;
  and an infinitesimally rigid framework attains the rank count exactly.
\end{definition}

\begin{theorem}[Realization at the target rank; KT Theorem 5.5]
  \label{thm:theorem-55}
  \lean{CombinatorialRigidity.Molecular.PanelHingeFramework.theorem_55_gen,
        CombinatorialRigidity.Molecular.PanelHingeFramework.theorem_55_minimalKDof_gen}
  \leanok
  \uses{def:rank-hypothesis, thm:minimal-kdof-reduction, lem:case-I,
        lem:case-II, lem:case-I-realization, lem:case-I-dispatch,
        lem:case-II-realization, lem:case-III, lem:theorem-55-base,
        lem:case-III-nested-rank-lower}
  Every minimal $0$-dof-graph $G$ with $|V| \ge 2$ has a panel-hinge
  realization $(G,p)$ with $\operatorname{rank} R(G,p) = D(|V(G)|-1)$ ---
  an infinitesimally rigid (on $V(G)$) panel-hinge framework
  (\cref{def:rank-hypothesis}, at $k' = 0$). Formalized at general $d$ for
  $D = \binom{n+1}{2} \ge 6$, i.e. $n \ge 3$ (the Phase-20 chain extractors
  fix this floor); a fresh-edge supply rides as ambient data. (The molecular
  conjecture needs only the $0$-dof instance stated here. The proof is not
  self-contained in $k$: Katoh--Tanigawa's induction runs over minimal
  $k$-dof-graphs for \emph{all} $k$, and the Case~III argument applies the
  inductive statement to an auxiliary minimal $k'$-dof graph with
  $0 \le k' \le D-2$ (Claim~6.11, eq.~(6.22)), discharged from the all-$k$ IH
  (\cref{lem:case-III-nested-rank-lower}).)
\end{theorem}
\begin{proof}
  \leanok
  \uses{thm:minimal-kdof-reduction, lem:theorem-55-base, lem:case-I,
        lem:case-II, lem:case-I-realization, lem:case-I-dispatch,
        lem:case-II-realization, lem:case-III, lem:case-III-nested-rank-lower,
        lem:chain-data-extract}
  Induction on $|V(G)|$, driven by the combinatorial reduction
  Theorem~4.9 (\cref{thm:minimal-kdof-reduction}): the all-$k$ well-founded
  induction principle (\texttt{minimal\_kdof\_reduction\_all\_k}) is run
  against the simplicity-conditioned motive
  ``$(G.\mathrm{Simple} \Rightarrow \mathrm{GP}) \wedge
    \mathrm{HasPanelRealization}$'' (\cref{def:rank-hypothesis}). The base
  case $|V(G)| = 2$ is \cref{lem:theorem-55-base}. The inductive step splits
  on: a cut edge (Case~0, \cref{lem:case-II-realization}); a proper rigid
  subgraph (Case~I, \cref{lem:case-I-dispatch}); $k > 0$ with no rigid
  subgraph (Case~II, \cref{lem:case-II-realization}); and $k = 0$ with no
  rigid subgraph (Case~III, \cref{lem:case-III}). The nested all-$k$ IH
  carries through all branches; Case~III's nested rank bound is
  \cref{lem:case-III-nested-rank-lower}. The grade-general spine
  \texttt{theorem\_55\_minimalKDof\_gen} fills the four reduction-case
  producers from their grade-general forms (base, cut, Case-I contract,
  and the forgetful map), with the Case-III arm run at general $n$ by the
  chain-dispatch router inside \texttt{case\_III\_realization\_all\_k} and
  the general chain extractor (\cref{lem:chain-data-extract}); the $d = 3$
  specialization is \cref{thm:theorem-55-d3-instance}.
\end{proof}

\begin{theorem}[Theorem 5.5 at $d = 3$: the zero-carry instance; KT Theorem 5.5, \S6.4.1]
  \label{thm:theorem-55-d3-instance}
  \lean{CombinatorialRigidity.Molecular.PanelHingeFramework.theorem_55_all_k,
        CombinatorialRigidity.Molecular.theorem_55_base_producer,
        CombinatorialRigidity.Molecular.PanelHingeFramework.theorem_55_d3,
        Graph.not_simple_of_isMinimalKDof_of_ncard_two,
        CombinatorialRigidity.Molecular.PanelHingeFramework.rankHypothesis_deficiency_of_theorem_55_d3}
  \leanok
  \uses{def:rank-hypothesis, def:genuine-hinge-realization,
        thm:minimal-kdof-reduction, lem:theorem-55-base-producer,
        lem:case-I-realization, lem:case-I-dispatch,
        lem:case-II-realization, lem:case-III, lem:case-III-nested-rank-lower}
  The $d = 3$ ($D = 6$) instance of \cref{thm:theorem-55}, run against the
  simplicity-conditioned motive: every minimal $0$-dof-graph $G$ with
  $|V| \ge 2$ satisfies the pair --- if $G$ is simple, a general-position,
  link-recording realization at an algebraically-independent seed of rank
  $D(|V|-1)$; and in all cases a genuine-hinge panel realization at rank
  $D(|V|-1)$ (\cref{def:genuine-hinge-realization}). The all-$k$ recursion
  (\texttt{minimal\_kdof\_reduction\_all\_k}) is run against the conditioned
  motive; all three adjudicated carries ($\mathtt{h622}$, $\mathtt{hsplit}$,
  $\mathtt{hcontract}$) are now discharged (Phase~22k~L7--L9): $\mathtt{h622}$
  via \cref{lem:case-III-nested-rank-lower}; $\mathtt{hsplit}$ by
  \cref{lem:simple-minimal-noRigid} (G0) $+$ Case-III
  (\cref{lem:case-III}); $\mathtt{hcontract}$ by \cref{lem:case-I-dispatch}.
  A fresh-edge supply rides as ambient data.

  \emph{The first push toward Theorem 5.6.} For a \emph{spanning} simple
  minimal $0$-dof-graph ($V(G)$ the whole body type), the general-position
  conjunct yields a panel-hinge framework realizing the rank hypothesis at
  $\operatorname{def}(\tilde G) = 0$ --- the $\operatorname{def} = 0$ generic-rank
  feed of the Proposition-1.1 reconciliation (\cref{prop:rigidity-matrix-prop11}).
  The $\operatorname{def} > 0$ form is the next theorem
  (\cref{thm:theorem-55-6-d3}).
\end{theorem}
\begin{proof}
  \leanok
  \uses{thm:minimal-kdof-reduction, lem:theorem-55-base-producer,
        lem:case-I-realization, lem:case-I-dispatch,
        lem:case-II-realization, lem:case-III, lem:case-III-nested-rank-lower,
        prop:rigidity-matrix-prop11}
  Instantiate \texttt{theorem\_55\_all\_k} at $D = 6$. The two-vertex base is
  \cref{lem:theorem-55-base-producer}. The no-rigid-subgraph branches
  ($k = 0$: \cref{lem:case-III}; $k > 0$: \cref{lem:case-II-realization}) each
  start with G0 to get $G.\mathrm{Simple}$, then apply the GP producer and
  M4 for the bare conjunct. The contraction branch dispatches on $k$:
  at $k = 0$, use \cref{lem:case-I-dispatch}; at $k > 0$, the 6.5 arm is
  vacuous (the carrier construction forces $k = 0$), and the remaining arms
  use \cref{lem:case-I-realization}. For the spanning stratum corollary,
  re-aim the endpoint selector at a fixed distinct pair on non-edges and apply
  the Proposition-1.1 rank count (\cref{prop:rigidity-matrix-prop11}).
\end{proof}

\begin{theorem}[Theorem 5.6 at $d = 3$: every multigraph attains the deficiency rank; KT Theorem 5.6, \S5.2]
  \label{thm:theorem-55-6-d3}
  \lean{CombinatorialRigidity.Molecular.PanelHingeFramework.rankHypothesis_of_theorem_55_d3,
        CombinatorialRigidity.Molecular.PanelHingeFramework.rankHypothesis_deficiency_of_theorem_55_d3}
  \leanok
  \uses{thm:theorem-55-d3-instance, lem:motions-mono-of-graph-le,
        lem:subgraph-minimality, def:rank-hypothesis}
  At $d = 3$ ($D = 6$), every simple spanning multigraph $G$ has a
  panel-hinge realization $(G,p)$ attaining the deficiency rank
  \[
    \operatorname{rank} R(G,p) = D(|V|-1) - \operatorname{def}(\tilde G),
  \]
  i.e.\ a framework realizing the rank hypothesis
  (\cref{def:rank-hypothesis}) at $k' = \operatorname{def}(\tilde G)$. This is
  Katoh--Tanigawa's Theorem~5.6 specialized to $d = 3$, lifting the
  minimal-$0$-dof realization of \cref{thm:theorem-55-d3-instance} to an
  arbitrary deficiency. It is the genuine $\operatorname{def} > 0$ feed of the
  Proposition-1.1 reconciliation (\cref{prop:rigidity-matrix-prop11}); the
  $\operatorname{def} = 0$ stratum is \cref{thm:theorem-55-d3-instance}.
\end{theorem}
\begin{proof}
  \leanok
  \uses{thm:theorem-55-d3-instance, lem:motions-mono-of-graph-le,
        lem:subgraph-minimality, prop:rigidity-matrix-prop11}
  Following Katoh--Tanigawa (p.~670). When $|V| = 1$ the framework is trivially
  rigid and $\operatorname{def}(\tilde G) = 0$, settled directly. Otherwise:
  delete edges of $G$ down to a minimal $k$-dof spanning subgraph $G' \le G$
  with $V(G') = V(G)$ and $\operatorname{def}(\tilde{G'}) =
  \operatorname{def}(\tilde G) =: k$ (the deficiency-preserving strip; the
  matroid engine behind \cref{lem:subgraph-minimality}). Realize $G'$ at its own
  deficiency $k$ by Theorem~5.5 (\cref{thm:theorem-55-d3-instance}, the all-$k$
  spine: $G'$ is simple as a subgraph of a simple $G$), yielding a generic
  full-rank realization $q$ on $G'$ with $\dim Z(G',q) = D + k$. Re-aim $q$ to
  the full graph $G$, keeping the supporting extensors on every $G'$-link
  unchanged and placing a genuine in-panel hinge on each re-added edge --- in
  the homogeneous model two distinct hyperplanes through the origin always meet,
  so KT's projective-invariance move is free. Re-adding edges only shrinks the
  motion space (\cref{lem:motions-mono-of-graph-le}), so
  $\dim Z(G,p) \le \dim Z(G',q) = D + k = D + \operatorname{def}(\tilde G)$. This
  is exactly the upper bound $\mathtt{hgen}$, which together with the
  genericity-free lower bound discharges the rank hypothesis via the
  Proposition-1.1 rank bridge (\cref{prop:rigidity-matrix-prop11}).
\end{proof}

\begin{lemma}[Theorem 5.5 base case ($|V(G)| = 2$)]
  \label{lem:theorem-55-base}
  \lean{CombinatorialRigidity.Molecular.BodyHingeFramework.theorem_55_base}
  \leanok
  \uses{def:rank-hypothesis, lem:rank-parallel-full}
  The two-vertex double edge has a panel-hinge realization of rank
  $D(|V(G)|-1) - k = D - k$. For the minimal $0$-dof case (the double edge,
  $k = 0$) this is the full rank $D$, attained by two non-parallel hinges
  via the parallel-hinges-full Lemma~5.3 (\cref{lem:rank-parallel-full}).
  At $k = 0$, in the $V(G)$-relative motive (\cref{def:rank-hypothesis}),
  the conclusion is exactly ``every infinitesimal motion is constant on the
  two bodies $V(G) = \{u, v\}$'' ($\mathrm{IsInfinitesimallyRigidOn}\,\{u,v\}$).
  The prior Lean required the constancy to hold over the \emph{whole} ambient
  body type (a hypothesis $\forall w,\ w = u \lor w = v$, i.e.\ ``$\alpha =
  \{u,v\}$''), the absolute-motive artefact; against the relativized motive the
  constancy is only required on $V(G)$, so that ambient hypothesis is
  dropped and the base case holds for any ambient type.
\end{lemma}
\begin{proof}
  \leanok
  \uses{lem:rank-parallel-full}
  Two bodies joined by hinges with linearly independent supporting
  extensors realize the full screw-difference constraint: the
  parallel-hinges-full Lemma~5.3 (\cref{lem:rank-parallel-full}) forces
  $S(u) = S(v)$ for every infinitesimal motion $S$. That is exactly
  constancy on $V(G) = \{u, v\}$, so the rank of the two hinge-row blocks is
  the full $D = D(2-1)$ and the $V(G)$-relative motive holds at $k = 0$. No
  condition on bodies outside $\{u, v\}$ is needed.
\end{proof}

\begin{lemma}[Two independent extensors in a common panel]
  \label{lem:extensor-pair-in-panel}
  \lean{CombinatorialRigidity.Molecular.exists_linearIndependent_extensor_pair_perp}
  \leanok
  \uses{def:extensor, lem:extensor-independence}
  For any normal $n \in \R^4$ there are two point-pairs $p, q : \{0,1\} \to \R^4$, each lying
  in the panel $n^{\perp} = \{x : x \cdot n = 0\}$, whose grade-2 extensors
  $\mathtt{extensor}\,p,\ \mathtt{extensor}\,q \in \bigwedge^2 \R^4$ are linearly independent.
  This is the coincident-panel ($\Pi(u) = \Pi(v) = n^{\perp}$) geometric brick of
  Katoh--Tanigawa's parallel-pair realization (Lemma~5.3, KT~2011 p.~670): two hinges sharing a
  panel but carrying non-proportional supporting extensors, the input the two-vertex base producer
  feeds to \cref{lem:theorem-55-base}.
\end{lemma}
\begin{proof}
  \leanok
  \uses{lem:extensor-independence}
  The panel $n^{\perp}$ has dimension at least $3$ in $\R^4$ (it is the kernel of the single
  pairing functional $x \mapsto x \cdot n$, so $\dim n^{\perp} \ge 4 - 1 = 3$, with equality
  unless $n = 0$). Pick three linearly independent vectors $a, b, c \in n^{\perp}$ and set
  $p = (a, b)$, $q = (a, c)$. The two wedges $a \wedge b$ and $a \wedge c$ are linearly
  independent: a relation $s\,(a \wedge b) + t\,(a \wedge c) = 0$, joined on the left with $c$
  (resp.\ $b$), kills the second (resp.\ first) term by the alternating property and leaves
  $s\,(c \wedge a \wedge b) = 0$ (resp.\ $t\,(b \wedge a \wedge c) = 0$); the surviving triple
  wedge is nonzero by independence of $\{a, b, c\}$, so $s = t = 0$. (This is the two-subset
  instance of the join-to-isolate technique behind \cref{lem:extensor-independence}.)
\end{proof}

\begin{lemma}[Theorem 5.5 base producer, empty arm]
  \label{lem:theorem-55-base-producer-empty}
  \lean{CombinatorialRigidity.Molecular.theorem_55_base_producer_empty,
        CombinatorialRigidity.Molecular.theorem_55_base_producer_empty_gp}
  \leanok
  \uses{def:rank-hypothesis, def:genuine-hinge-realization, lem:two-vertex-trichotomy}
  Let $G$ be a minimal-$k$-dof-graph on $1 \le |V(G)| \le 2$ bodies realizing the \emph{empty}
  case of the base trichotomy (\cref{lem:two-vertex-trichotomy} (i)): $E(G) = \emptyset$ and
  hence $\mathrm{def}(\tilde G) = D(|V(G)| - 1)$. Then $G$ has a genuine-hinge panel realization
  (\cref{def:genuine-hinge-realization}) at the target rank $D(|V(G)|-1) - \mathrm{def} = 0$, and
  --- the general-position companion --- a \emph{generic} full-rank realization
  (\cref{def:rank-hypothesis}) at that same rank.
\end{lemma}
\begin{proof}
  \leanok
  The bare construction is geometry-free: take the all-zero framework (every edge carries the
  zero supporting extensor) at a fixed nonzero panel normal. With no edges no hinge constraint
  fires, so the rigidity-row family is empty, its span is $\{0\}$, and the rank is $0$; the
  per-link conjunct of the motive is vacuous. The target rank vanishes because the empty-graph
  deficiency is exactly $D(|V(G)|-1)$ (\cref{lem:two-vertex-trichotomy}). For the generic
  companion build the single-body / empty framework at an injective
  algebraically-independent seed (a non-root of the general-position polynomial): the panel
  normals are then in general position and algebraically independent, the rigidity-row family is
  still empty (rank $0$), and link-recording is vacuous.
\end{proof}

\begin{lemma}[Theorem 5.5 base producer, single-edge arm]
  \label{lem:theorem-55-base-producer-single}
  \lean{CombinatorialRigidity.Molecular.theorem_55_base_producer_single_edge,
        CombinatorialRigidity.Molecular.theorem_55_base_producer_single_edge_gp}
  \leanok
  \uses{def:rank-hypothesis, def:genuine-hinge-realization, lem:two-vertex-trichotomy,
        lem:extensor-pair-in-panel, def:rigidity-matrix}
  Let $G$ be a two-vertex minimal-$1$-dof-graph realizing the \emph{single-edge} case of the base
  trichotomy (\cref{lem:two-vertex-trichotomy} (ii)): $V(G) = \{x, y\}$ with $x \ne y$,
  $E(G) = \{e\}$ a single edge joining $x$ and $y$, and $\mathrm{def}(\tilde G) = 1$. Then $G$ has
  a genuine-hinge panel realization (\cref{def:genuine-hinge-realization}) at the target rank
  $D(|V(G)|-1) - \mathrm{def} = D - 1 = 5$ (at $d = 3$, $D = 6$), and --- the general-position
  companion, the one base arm where this conjunct does genuine geometric work --- a \emph{generic}
  full-rank realization (\cref{def:rank-hypothesis}) at that same rank.
\end{lemma}
\begin{proof}
  \leanok
  \uses{lem:extensor-pair-in-panel, def:rigidity-matrix}
  For the bare motive, place coincident panels $\Pi(x) = \Pi(y) = n^{\perp}$ at a fixed nonzero
  normal $n$ and assign the single hinge one nonzero supporting extensor inside the common panel
  (the first of the pair of \cref{lem:extensor-pair-in-panel}; nonzero by its linear
  independence). The lone hinge's row block already \emph{is} the entire rigidity-row family
  (there is only one edge), and that block has rank exactly $D - 1$ (\cref{def:rigidity-matrix});
  so the rigidity-row span has rank $D - 1 = D(|V(G)|-1) - \mathrm{def}$, the rank equality of the
  motive. Unlike the parallel-pair arm the framework is not rigid ($\mathrm{def} = 1$): the rank
  is the single hinge-row count, not the full two-vertex base rank. The generic companion builds
  $\mathtt{ofNormals}$ at an injective algebraically-independent seed (a non-root of the
  general-position polynomial) with the single edge's two distinct endpoints: general position
  forces the lone supporting extensor nonzero, the same single hinge-row count gives rank $D - 1$,
  and the seed supplies general position $+$ algebraic independence of the panel normals.
\end{proof}

\begin{lemma}[Theorem 5.5 base producer, parallel-pair arm]
  \label{lem:theorem-55-base-producer-parallel}
  \lean{CombinatorialRigidity.Molecular.theorem_55_base_producer_parallel_pair}
  \leanok
  \uses{def:genuine-hinge-realization, lem:two-vertex-trichotomy, lem:theorem-55-base,
        lem:extensor-pair-in-panel}
  Let $G$ be a two-vertex minimal-$0$-dof-graph realizing the \emph{parallel-pair} case of the
  base trichotomy (\cref{lem:two-vertex-trichotomy} (iii)): $V(G) = \{x, y\}$ with $x \ne y$,
  $E(G) = \{e, f\}$ with $e \ne f$ two edges both joining $x$ and $y$, and $\mathrm{def}(\tilde G)
  = 0$. Then $G$ has a genuine-hinge panel realization
  (\cref{def:genuine-hinge-realization}) at the full target rank $D(|V(G)|-1) - \mathrm{def} =
  D \cdot 1 = 6$ (at $d = 3$, $D = 6$). This is the geometrically nontrivial arm of the all-$k$
  base producer; the empty and single-edge arms are bookkeeping / single-row counts.
\end{lemma}
\begin{proof}
  \leanok
  \uses{lem:theorem-55-base, lem:extensor-pair-in-panel}
  Place \emph{coincident panels} $\Pi(x) = \Pi(y) = n^{\perp}$ at a fixed nonzero normal
  $n$ and assign the two parallel hinges two linearly independent supporting extensors inside the
  common panel $n^{\perp}$ (\cref{lem:extensor-pair-in-panel}). Each link's supporting extensor is
  then nonzero and lies in both endpoint panels, discharging the per-link conjunct of the motive.
  The two independent extensors give the combined hinge-row blocks full rank $D$ on the relative
  screw $S(x) - S(y)$, so the two-vertex base case (\cref{lem:theorem-55-base}) makes the framework
  infinitesimally rigid on $\{x, y\} = V(G)$; the $V(G)$-relative rank bridge (rigid on $V(G)$ iff
  the rigidity-row span has rank $D(|V(G)|-1)$) turns that into the rank equality of the motive,
  with $\mathrm{def}(\tilde G) = 0$.
\end{proof}

\begin{lemma}[Theorem 5.5 base producer, trichotomy dispatch]
  \label{lem:theorem-55-base-producer}
  \lean{CombinatorialRigidity.Molecular.theorem_55_base_producer}
  \leanok
  \uses{def:rank-hypothesis, def:genuine-hinge-realization, lem:two-vertex-trichotomy,
        lem:theorem-55-base-producer-empty, lem:theorem-55-base-producer-single,
        lem:theorem-55-base-producer-parallel}
  Let $G$ be a minimal-$k$-dof-graph with $1 \le |V(G)| \le 2$ (the base region of the
  all-$k$ reduction, \cref{thm:minimal-kdof-reduction-all-k}).  Then the conditioned pair
  $(G\text{ simple} \Rightarrow \text{has a \emph{generic} full-rank panel realization})$
  and $(\text{has a genuine-hinge panel realization at target rank})$ both hold.  The
  simple conjunct is the strengthened generic motive
  (\cref{def:rank-hypothesis}, general position $+$ algebraically-independent panel
  normals); the unconditional conjunct is the honest bare motive
  (\cref{def:genuine-hinge-realization}).  This is the conditioned pair $P_c\,G$ that the
  all-$k$ spine recurses against.
\end{lemma}
\begin{proof}
  \leanok
  \uses{lem:theorem-55-base-producer-empty, lem:theorem-55-base-producer-single,
        lem:theorem-55-base-producer-parallel}
  Dispatch on the base trichotomy (\cref{lem:two-vertex-trichotomy}). The unconditional
  bare conjunct comes from the three bare arms. For the simple conjunct: the empty arm
  ($E(G) = \emptyset$) and the single-edge arm ($\mathrm{def} = 1 > 0$) both admit a
  \emph{generic} full-rank realization --- built from $\mathtt{ofNormals}$ at an injective
  algebraically-independent seed, which is a non-root of the general-position polynomial, so
  the panel normals are simultaneously in general position and algebraically independent; the
  rank is $0$ (empty) resp.\ $D - 1$ (single-edge, the single hinge-row count). The
  single-edge arm is the one base arm where the simple conjunct does genuine geometric work.
  In the parallel-pair arm ($k = 0$) a two-vertex minimal-$0$-dof-graph cannot be simple (the
  single-vertex cut supplies two distinct parallel edges, both forced to link the same pair,
  which simplicity collapses --- contradiction); the $G$-simple antecedent is therefore
  vacuous.
\end{proof}

\begin{lemma}[Theorem 5.5 triangle base ($|V(G)| = 3$)]
  \label{lem:theorem-55-triangle}
  \lean{CombinatorialRigidity.Molecular.BodyHingeFramework.theorem_55_triangle}
  \leanok
  \uses{def:rank-hypothesis, lem:rank-parallel-full}
  The three-body triangle sibling of the two-body base case
  \cref{lem:theorem-55-base}: if a body-hinge framework $F$ has three
  edges $e_1 : u$--$v$, $e_2 : v$--$w$, $e_3 : w$--$u$ forming a
  triangle whose three supporting extensors $C(p(e_1))$, $C(p(e_2))$,
  $C(p(e_3))$ are linearly independent, then $F$ is infinitesimally rigid
  on the three bodies $\{u, v, w\}$
  ($\mathrm{IsInfinitesimallyRigidOn}\,\{u,v,w\}$).
\end{lemma}
\begin{proof}
  \leanok
  \uses{lem:rank-parallel-full}
  Each hinge constraint puts a relative screw $S(u) - S(v)$, $S(v) -
  S(w)$, $S(w) - S(u)$ in the one-dimensional span of its supporting
  extensor; the three differences telescope to $0$ (the cyclic shift is a
  bijection of $\{u,v,w\}$). Linear independence of the three extensors
  --- the $m = 3$ instance of \cref{lem:rank-parallel-full}'s
  $m$-dimensional generalization --- forces each difference to vanish, so
  $S(u) = S(v) = S(w)$, and constancy on $\{u,v,w\}$ follows.
\end{proof}

\begin{lemma}[Cyclic-seed existence for the triangle base]
  \label{lem:triangle-normals}
  \lean{CombinatorialRigidity.Molecular.exists_triangle_normals}
  \leanok
  \uses{def:panel-support-extensor, lem:panel-support-extensor-independence,
        lem:extensor-independence}
  When $k \ge 1$ (so $k+2 \ge 3$), there exist three vectors $n_0, n_1, n_2 \in \R^{k+2}$
  such that (1) each cyclic pair $(n_0,n_1)$, $(n_1,n_2)$, $(n_2,n_0)$ has a nonzero
  grade-2 join ($n_i \wedge n_j \ne 0$), and (2) the cyclic supporting-extensor family
  $i \mapsto \mathrm{panelSupportExtensor}(n_{\sigma(i)}, n_{\sigma(i)+1})$ is linearly
  independent. The witness is $n_0 = e_0$, $n_1 = e_1$, $n_2 = e_2$ (standard basis
  vectors in $\R^{k+2}$): the cyclic family reduces (via the grade-2 antisymmetry at the
  reversed pair $(e_2, e_0) = -(e_0 \wedge e_2)$) to the sorted family $e_0 \wedge e_1$,
  $e_1 \wedge e_2$, $e_0 \wedge e_2$, up to unit scalars; the sorted family is a
  three-member subfamily of the $\Lambda^2$-basis indexed by distinct 2-subsets; that
  basis family is linearly independent (\cref{lem:extensor-independence}); unit scaling
  preserves linear independence; and $e_i \wedge e_j \ne 0$ for $i < j$ since each such
  join is a nonzero member of the same linearly independent basis family.
\end{lemma}
\begin{proof}
  \leanok
  \uses{lem:panel-support-extensor-independence, lem:extensor-independence}
  By \cref{lem:panel-support-extensor-independence} it suffices to show the grade-2 join
  family is linearly independent. Via \texttt{normalsJoin\_swap}, the cyclic join family
  $(e_0 \wedge e_1,\, e_1 \wedge e_2,\, e_2 \wedge e_0)$ equals (unit scalars) times the
  sorted family $(e_0 \wedge e_1,\, e_1 \wedge e_2,\, e_0 \wedge e_2)$; unit scaling
  preserves independence. Each sorted join equals the corresponding member of the
  $\Lambda^2$-basis family (\texttt{normalsJoin\_basisFun\_orderEmbOfFin}); the three
  indices $\{0,1\}$, $\{1,2\}$, $\{0,2\}$ are distinct, so the three-element subfamily
  inherits independence from the full basis family
  (\cref{lem:extensor-independence}).
\end{proof}

\begin{lemma}[Cyclic-seed existence for a general cycle]
  \label{lem:cycle-normals}
  \lean{CombinatorialRigidity.Molecular.exists_cycle_normals}
  \leanok
  \uses{def:panel-support-extensor, lem:panel-support-extensor-independence,
        lem:extensor-independence}
  For $3 \le m \le k + 2$ there exist $m$ vectors $n_0, \dots, n_{m-1} \in \R^{k+2}$
  such that (1) each cyclic pair $(n_i, n_{i+1})$ (indices mod $m$) has a nonzero
  grade-2 join $n_i \wedge n_{i+1} \ne 0$, and (2) the cyclic supporting-extensor
  family $i \mapsto \mathrm{panelSupportExtensor}(n_i, n_{i+1})$ is linearly
  independent. This is the general-$m$ generalization of the triangle seed
  \cref{lem:triangle-normals} (its $m = 3$ instance), and the \emph{shared}-normal
  sibling of the free-pair extensor existence
  \cref{lem:exists-independent-panel-extensor}: consecutive hinges of a panel cycle
  share a panel, so a panel-hinge seed assigns \emph{one} normal per body --- the
  free pairs of \cref{lem:exists-independent-panel-extensor} cannot feed it.
\end{lemma}
\begin{proof}
  \leanok
  \uses{lem:panel-support-extensor-independence, lem:extensor-independence}
  Witness: the standard basis restricted along the inclusion of index sets,
  $n_i = e_i$ (this is where $m \le k + 2$ enters). The $m$ cyclic joins then
  realize the $m$ distinct $2$-subsets $\{i, i+1\}$ ($i < m-1$) and $\{0, m-1\}$ ---
  the edges of the cycle $C_m$ embedded in the index set. Via the grade-2
  antisymmetry at the single wrap edge, the cyclic join family equals unit scalars
  times the sorted subfamily of the $\Lambda^2$-basis family indexed by these
  $2$-subsets; the edge-to-subset index map is injective for $m \ge 3$ (two cyclic
  edges coincide only if $2 = 0$ in $\Z/m$), so the subfamily inherits independence
  from the basis family (\cref{lem:extensor-independence}), and unit scaling
  preserves independence. Each join is nonzero since consecutive indices are
  distinct; the extensor conclusion is carried across
  \cref{lem:panel-support-extensor-independence}.
\end{proof}

\subsection{The framework-construction op}
\label{sec:molecular-algebraic-induction-withgraph}

Both inductive cases build a realization of $G$ from a realization of a
\emph{different} graph: Case~I from the contraction $G/E(H)$, Case~II
from the splitting-off $G_v^{ab}$. The shared, citation-free primitive
both need is the ability to keep the hinge assignment fixed while
changing the underlying multigraph --- the supporting extensors $C(p(e))$
and hinge-row blocks depend only on the hinge data, not on which edges
$G$ carries. The one fact this phase needs from the op is the
graph-monotonicity of the motion space: deleting edges (passing to a
subgraph) can only enlarge the null space $Z(G,p)$, the direction
Cases~I and~II travel toward the smaller inductive graph.

\begin{definition}[Framework on a new graph]
  \label{def:framework-with-graph}
  \lean{CombinatorialRigidity.Molecular.BodyHingeFramework.withGraph,
        CombinatorialRigidity.Molecular.PanelHingeFramework.withGraph,
        CombinatorialRigidity.Molecular.PanelHingeFramework.toBodyHinge_withGraph,
        CombinatorialRigidity.Molecular.PanelHingeFramework.withNormal,
        CombinatorialRigidity.Molecular.PanelHingeFramework.toBodyHinge_withNormal_supportExtensor_of_ne,
        CombinatorialRigidity.Molecular.PanelHingeFramework.toBodyHinge_withNormal_infinitesimalMotions_eq,
        CombinatorialRigidity.Molecular.PanelHingeFramework.toBodyHinge_withNormal_pinnedMotions_eq}
  \leanok
  \uses{def:rigidity-matrix, def:panel-hinge-framework}
  For a body-hinge framework $F = (G, p)$ and any multigraph $G'$ on the
  same bodies, $F[G']$ is the framework with underlying multigraph $G'$
  and the same hinge assignment $p$. Every supporting extensor $C(p(e))$,
  hinge-row block, and per-edge constraint is untouched; only the set of
  edges imposing a constraint changes. The panel layer carries the same
  operation on a panel-hinge framework $P$ --- swap the multigraph, keep
  the per-body panel normals and endpoint selector --- and the two commute
  with the body-hinge interpretation, $(P[G'])^\flat = (P^\flat)[G']$, so
  the graph-monotonicity and block-pin rank facts below apply verbatim to a
  panel realization re-glued onto $G$, with coplanarity preserved (the
  normals are unchanged). Case~II's $1$-extension additionally re-inserts a
  body $v$ with a \emph{chosen} panel: the per-body variant $P[v \mapsto n]$
  overrides only $v$'s panel normal, leaving the multigraph, endpoint
  selector, and every other normal fixed, so each hinge whose endpoints
  avoid $v$ keeps its supporting extensor and the inductive realization of
  the splitting-off $G_v^{ab}$ is untouched away from $v$. In particular,
  when $v$ is still unhinged (no edge of $G$ has $v$ as an endpoint, as in
  $G_v^{ab}$ before $v$'s two new edges are added), the choice of $v$'s
  panel normal changes neither the null space $Z(G,p)$ nor any
  body-$w$-pinned motion subspace, so the inductive count is preserved
  through the choice.
\end{definition}

\begin{lemma}[Edge deletion enlarges the motion space]
  \label{lem:motions-mono-of-graph-le}
  \lean{CombinatorialRigidity.Molecular.BodyHingeFramework.infinitesimalMotions_le_withGraph_of_le,
        CombinatorialRigidity.Molecular.BodyHingeFramework.finrank_infinitesimalMotions_le_of_graph_le}
  \leanok
  \uses{def:framework-with-graph}
  If $G' \le G$ then $Z(G,p) \subseteq Z(G', p)$, equivalently
  $\dim Z(G,p) \le \dim Z(G', p)$ --- i.e.
  $\operatorname{rank} R(G', p) \le \operatorname{rank} R(G, p)$. A
  supergraph imposes more hinge constraints, so a motion of $G$ is a
  fortiori a motion of any subgraph $G'$. This is the
  ``re-adding edges only grows the rank'' monotonicity that lifts a
  realization of a minimal $k$-dof spanning subgraph to one of the whole
  multigraph (the step \cref{prop:rigidity-matrix-prop11} uses), and the
  companion to the fixed-graph span-monotonicity Lemma~5.2.
\end{lemma}
\begin{proof}
  \leanok
  \uses{def:framework-with-graph}
  Each link of $G'$ is a link of $G$ (\texttt{Graph.IsLink.mono}), and
  the matching supporting extensors agree, so every infinitesimal motion
  of $G$ meets every constraint of $G'$: $Z(G,p) \subseteq Z(G', p)$.
  Taking dimensions (\texttt{Submodule.finrank\_mono}) gives the rank
  form.
\end{proof}

